%% VI. Discussion: interpretation by research question, implications, limitations. Input by access.tex.
\section{Discussion}
\label{sec:discussion}

\subsection{What the Design Analysis Got Right (RQ1)}
The companion paper promised that the exact Fisher information of M1 could be used to plan an experience-sampling study. On two real designs, with irregular probes, a third of the epochs missing and sessions excluded by an artefact rule, the bound computed before the lock, the Fisher standard error at the estimate and the bootstrap error of $\kinh$ differed by at most 8\,\% on ds001787 and 32\,\% on OSF 3uszv. On ds001787 the achieved 9.9\,\% lay between the bound computed before the lock and the bound recomputed at the measured signal-to-noise ratio (9.1--13.3\,\%); on OSF 3uszv the achieved 8.5\,\% exceeded the recomputed bound (7.4\,\%) by 15\,\%. For $\kd$ the bound computed before the lock was optimistic by a factor of 1.8 and 1.2. The ordering between the datasets was also predicted: the denser probing and longer sessions of OSF 3uszv give a smaller bound and a smaller achieved error. A design analysis of this kind can therefore be used for its intended purpose, deciding how many subjects, how long a session and how many probes are needed, provided the error it quotes is understood as a Fisher error of a model assumed to be right.

\subsection{What It Cannot Anticipate (RQ1)}
A bound cannot anticipate misspecification, and the alpha-power channel is misspecified on both datasets: the observed-information and profile errors of $\kinh$, $\sA$ and $\rA$ exceed the Fisher error by a factor of 2.1--2.5 on ds001787 and 1.1--1.4 on OSF 3uszv, and the excess disappears on data simulated from the fit. The diagnostic is cheap, since it needs one Hessian and one simulated dataset, and we recommend it alongside the bootstrap, which cannot see misspecification because it simulates from the model. In practice, design-map numbers should be read as lower bounds whenever a proxy is not known to follow the model, and likelihood-based intervals should be reported beside Fisher-based ones; on ds001787 the two criteria for H1 disagree.

\subsection{The Observation Model Is the Bottleneck (RQ2)}
The two proxies we preregistered carried no detectable information about the reports beyond the reports' own serial dependence, on either dataset. The complexity channel was noise on ds001787 and, on OSF 3uszv, a 40-s signal unrelated to the reports; posterior alpha did not load on the ratings on either. The dynamics of M1 are identifiable and were estimated with the predicted precision, but the read-out from EEG features to the reported state is where the information is missing. This is consistent with the indirect relation between posterior alpha and the DMN found in simultaneous EEG--fMRI \cite{Laufs2003,Mantini2007}. Finding an EEG feature that tracks the reported state, with a signal-to-noise ratio in the range that the design analysis shows to be necessary, is the open problem that this class of models leaves to experimenters, and the experience-sampling datasets now available \cite{Weinstein2018,Brandmeyer2018,RodriguezLarios2021} are the place to look for it.

\subsection{Implications for the BQI Framework (RQ3)}
Neither dataset supports the predictions T1, T2, T4 and T5 of BQI, and T3 holds in sign only (Section~\ref{sec:res-bqi}), but the comparison is not like for like. The BQI values were fitted to meta-analytic summaries of fMRI contrasts and cross-sectional correlations \cite[Sec.~7.2--7.5]{BQI}, a limitation that BQI itself states \cite[Sec.~11.3]{BQI}, whereas the present estimates come from EEG proxies observed over minutes within single sessions. A time constant fitted to meta-analytic group summaries describes at most how group-level DMN contrasts evolve over a practice; the 0.76 and 7.1~min estimated here describe how posterior alpha power fluctuates around its session mean. Because the Ornstein--Uhlenbeck DMN of M1 is the quadratic-potential case of the BQI Langevin dynamics \cite[eqs.~(83)--(87)]{BQI}, the estimated $\kinh$ is the curvature of that potential for the EEG proxy, not for the DMN itself. The negative coupling on OSF 3uszv says that, with posterior alpha standing in for the DMN, alpha and complexity rise together, the opposite of the coupling in \cite[eq.~(53)]{BQI}; it does not refute that eq.~(53) holds for the DMN itself. What the comparison shows is the step that BQI leaves open: before its parameters can be estimated from EEG, the EEG features must be calibrated against the state variables, for instance with the simultaneous EEG, fMRI and complexity recordings that BQI itself proposes \cite[Sec.~14]{BQI}, and the QoC must be operationalised with a validated instrument. The structural parts of BQI that M1 keeps, an Ornstein--Uhlenbeck DMN process feeding an integration variable and a linear read-out to the report, were estimable on both datasets with the precision the design analysis predicted.

\subsection{Rate Constants Are Properties of the Proxy (RQ4)}
The replication separates two kinds of result. The methodological ones held across laboratories: predicted and achieved errors of $\kinh$ agreed to within a third, the alpha channel showed the same signature, the EEG added nothing to prediction, and the heterogeneity between practitioners was large. The substantive ones did not: the DMN time constant moved from 0.8 to 7~min, the integration time constant from 5 to 0.6~min, and the group difference in $\kinh$ that ds001787 had shown disappeared. The simulation of Section~\ref{sec:simmethods} shows that this is expected rather than puzzling. An Ornstein--Uhlenbeck fit returns a variance-weighted compromise between the time scales present in the feature; posterior alpha power in a 45-min mantra session and in a 60-min breath-focus session, recorded with different montages, need not contain the same mixture. The lesson is not that one of the two estimates is wrong but that a rate constant estimated from an uncalibrated proxy is a property of the proxy in that recording. The parameters that a design analysis can promise to deliver are only as meaningful as the mapping from features to state variables, and that mapping must be established and shown to be stable across laboratories before a design map is used to plan a study or a closed-loop intervention \cite{Sitaram2017,Ros2020}: the control gains computed in the companion paper assume parameters that mean something.

\subsection{What the Preregistration Caught}
Two episodes illustrate the value of locks and gates. The locked artefact rule of ds001787 excluded 45\,\% of the sessions, which the stop criterion allowed; without the lock the natural reaction would have been to switch to the adaptive rule, which gives a different integration time constant and a positive $\beta$, and to report that. On OSF 3uszv a sensitivity variant copied from the first preregistration could not run, and 17 sessions were first preprocessed on a machine without the ICLabel package, which the preprocessing logs recorded; they were reprocessed before any fit was reported, and both facts are listed in Appendix~\ref{app:deviations}. Preregistration did not make the analysis simpler, but it made the list of what was decided after seeing the data short and explicit \cite{Nosek2018}.

\subsection{Limitations}
\label{sec:limitations}
The empirical tests are two datasets whose rate constants disagree by factors of eight and nine. On ds001787 the preregistered artefact rule excluded 18 of 40 sessions non-randomly, and the group comparison is confounded by sex and age; on OSF 3uszv the probe windows alone remove a third of the EEG, the ratings are bimodal and serially independent, and the sessions include an initial segment before the first probe whose masking doubles the DMN time constant. On both, the EEG features are proxies without a calibrated mapping to $A$ or $\Phi$ (posterior alpha is not a measure of DMN activity, Lempel--Ziv complexity is not a measure of $\Phi$), the levels $A_0$, $\Phi^*$ and $\gamma_0$ are not estimable after centring, and the alpha channel shows signs of misspecification. The two rating scales are operational stand-ins for the QoC of BQI rather than validated measures of it, and the complexity proxy, like the index of BQI, measures complexity rather than integration \cite[Sec.~11.3]{BQI}. M1 is linear with Gaussian noise and a linear read-out, whereas the ratings are ordinal and bounded, bimodal on one dataset, and meditation dynamics may involve discrete state switches \cite{Hasenkamp2012,Christoff2016}; non-linear extensions would require approximate filters. Subjects share one parameter vector in the pooled fits, although the between-subject standard deviation of $\log\kinh$ is 1.2--1.5 log units; a hierarchical model is the natural next step and would change the design map quantitatively. Finally, the analysis was planned, locked and run by a single author, and the preregistrations were not reviewed before the lock.
